\chapter{Differentiating under the integral sign}
\label{app:feynman}

Exercise~\ref{ex:ecos} needs $\E[\cos\theta]$ for a Gaussian angle
$\theta$, and its solution finds it by a device that is hard to invent
on the spot: put a parameter $a$ into the integral, differentiate with
respect to $a$, and solve the equation that comes out. The device is
classical. Richard Feynman made it popular under the name
\emph{differentiating under the integral sign}; he had learned it as a
student from Woods' \emph{Advanced Calculus}, and it is now often called
\emph{Feynman's trick}. This appendix states the rule that makes it legal
(\cref{sec:feyn-rule}), describes the device (\cref{sec:feyn-idea}),
works three examples of the two ways it closes
(\cref{sec:feyn-easy,sec:feyn-ode}), gives the shorter route to
$\E[\cos\theta]$ through complex numbers (\cref{sec:feyn-complex}), and
ends with how to choose the parameter (\cref{sec:feyn-choose}).

\section{The rule}
\label{sec:feyn-rule}

All the integrals below are integrals of integrable functions, that is,
Lebesgue integrals, and expectations are integrals of this kind. The
rule is a consequence of dominated convergence.

\begin{proposition}{Differentiation under the integral sign}{diffint}
Let $J$ be an open interval, $S$ a subset of $\R$, and
$f:J\times S\to\R$ a function such that
\begin{enumerate}[nosep]
  \item for each $a\in J$, the function $x\mapsto f(a,x)$ is integrable
        on $S$;
  \item for each $x\in S$, the function $a\mapsto f(a,x)$ is
        differentiable on $J$;
  \item there is an integrable function $g$ on $S$ with
        $\abs{\partial_af(a,x)}\le g(x)$ for all $a\in J$ and $x\in S$.
\end{enumerate}
Then $F(a)=\int_Sf(a,x)\,\mathrm dx$ is differentiable on $J$, and
\[
  F'(a)=\int_S\partial_af(a,x)\,\mathrm dx .
\]
\end{proposition}

\begin{proof}
Fix $a\in J$ and a sequence $h_n\to0$ with $h_n\neq0$ and
$a+h_n\in J$. The difference quotient of $F$ is the integral of the
difference quotients of $f$:
\[
  \frac{F(a+h_n)-F(a)}{h_n}=\int_Sq_n(x)\,\mathrm dx,
  \qquad
  q_n(x)=\frac{f(a+h_n,x)-f(a,x)}{h_n}.
\]
For each $x$, $q_n(x)\to\partial_af(a,x)$ by (2), so $\partial_af(a,\cdot)$
is measurable as a limit of measurable functions. By the mean value
theorem, $q_n(x)=\partial_af(\xi_n,x)$ for some $\xi_n$ between $a$ and
$a+h_n$, so $\abs{q_n(x)}\le g(x)$ by (3). Dominated convergence then
gives $\int_Sq_n\to\int_S\partial_af(a,x)\,\mathrm dx$. Write $L=\int_S\partial_af(a,x)\,\mathrm dx$. We have shown: for every
sequence $h_n\to0$ with $h_n\neq0$ and $a+h_n\in J$, the difference
quotient $\bigl(F(a+h_n)-F(a)\bigr)/h_n$ tends to $L$.

It remains to pass from sequences to the limit as $h\to0$, which is what
differentiability asks for; sequences were needed because dominated
convergence is a statement about sequences. Suppose that the quotient $\bigl(F(a+h)-F(a)\bigr)/h$ did not tend to $L$
as $h\to0$. Then there would be an $\varepsilon>0$ such
that, however small $\delta>0$, some $h$ with $0<\abs h<\delta$ and
$a+h\in J$ gives $\bigl|\bigl(F(a+h)-F(a)\bigr)/h-L\bigr|\ge\varepsilon$.
Taking $\delta=1/n$ for $n=1,2,\dots$ gives such an $h_n$ for each $n$:
a sequence with $h_n\to0$, $h_n\neq0$, $a+h_n\in J$, whose quotients all
stay at distance at least $\varepsilon$ from $L$, and so do not tend to
$L$. This contradicts what was shown. So the quotient tends to $L$ as
$h\to0$: $F$ is differentiable at $a$, and $F'(a)=L$, the stated
derivative. Since $a$ was any point of $J$, this holds on all of $J$.
\end{proof}

A complex-valued $f$ is handled by applying the proposition to its real
and imaginary parts. The companion book proves a more general form, with
a parameter in a normed space and a Lipschitz bound in place of (3)
(Theorem~17.20 there).

\emph{Hypothesis (3) is the one to check.} Hypotheses (1) and (2) are
about one $a$, or one $x$, at a time, and they are usually clear.
Hypothesis (3) asks for more: a single integrable function $g$ that bounds
$\partial_af(a,\cdot)$ for \emph{all} $a\in J$ at once. It forbids the
derivative from carrying its mass off to infinity as $a$ varies, and that
can happen.

\emph{An example where (3) fails.} Take $J=(-1,1)$, $S=(0,\infty)$ and
\[
  f(a,x)=a^3x\,e^{-a^2x^2}.
\]
Hypotheses (1) and (2) hold: for each $a$ the function of $x$ is
integrable, and for each $x$ the function of $a$ is differentiable. For
$a\neq0$,
\[
  F(a)=a^3\int_0^\infty x\,e^{-a^2x^2}\,\mathrm dx
  =a^3\Bigl[-\frac{e^{-a^2x^2}}{2a^2}\Bigr]_0^\infty
  =a^3\cdot\frac1{2a^2}=\frac a2 ,
\]
and $F(0)=0$, so $F(a)=a/2$ on all of $J$, and $F'(0)=\tfrac12$. But
\[
  \partial_af(a,x)=x\,e^{-a^2x^2}\bigl(3a^2-2a^4x^2\bigr)
\]
is $0$ at $a=0$, for every $x$, so $\int_S\partial_af(0,x)\,\mathrm dx=0$:
the formula of \cref{prop:diffint} gives $0$ instead of $\tfrac12$.
Hypothesis (3) is what fails. For $x>1$, the value $a=1/x$ is in $J$,
and there $\partial_af=x\,e^{-1}(3/x^2-2/x^2)=e^{-1}/x$; so any $g$ that
bounds $\abs{\partial_af(a,x)}$ for all $a\in J$ has $g(x)\ge e^{-1}/x$
for $x>1$, and $\int_1^\infty\mathrm dx/x=\infty$: no such $g$ is
integrable. The picture: for small $a$, the derivative is a bump around
$x\approx1/a$, which runs off to infinity as $a\to0$, and its mass goes
with it.

\emph{An integral that is not one: $\int_0^\infty\sin(ax)/x\,\mathrm dx$.}
Hypothesis (1) can fail too. The function $\sin x/x$ is not integrable on
$(0,\infty)$: on each interval $[k\pi,(k+1)\pi]$ the function $\abs{\sin
x}$ has the integral $2$ and $1/x\ge1/((k+1)\pi)$, so
\[
  \int_0^\infty\frac{\abs{\sin x}}x\,\mathrm dx
  \ge\sum_{k=0}^\infty\frac2{(k+1)\pi}=\infty,
\]
like the harmonic series. So the integral of $\sin x/x$ over $(0,\infty)$,
in the sense used in this appendix, does not exist. What does exist is
the limit of the integrals over $(0,R)$ as $R\to\infty$: the positive and
negative arches cancel more and more, and
$\int_0^R\sin x/x\,\mathrm dx$ tends to a limit (it is $\pi/2$, though
the value is not needed here). Such a limit of integrals over growing
intervals is called an \emph{improper integral}; it is not the integral of
an integrable function, and \cref{prop:diffint} says nothing about it.

And the formula does fail. With $a>0$ and
$F(a)=\lim_{R\to\infty}\int_0^R\sin(ax)/x\,\mathrm dx$, the substitution
$y=ax$ turns $\int_0^R\sin(ax)/x\,\mathrm dx$ into
$\int_0^{aR}\sin y/y\,\mathrm dy$, which has the same limit as
$R\to\infty$; so $F(a)=F(1)$ for every $a>0$, and $F'(a)=0$.
Differentiating under the integral sign anyway would give
$\int_0^\infty\cos(ax)\,\mathrm dx$, which does not exist even as an
improper integral: $\int_0^R\cos(ax)\,\mathrm dx=\sin(aR)/a$ keeps
oscillating as $R\to\infty$.

\section{The device}
\label{sec:feyn-idea}

To compute one number $I$ given as an integral:
\begin{enumerate}[nosep]
  \item \emph{Embed it in a family.} Write $I=F(a_1)$ for a function
        $F(a)=\int f(a,x)\,\mathrm dx$ of a parameter $a$.
  \item \emph{Choose the family so that one member is known}:
        $F(a_0)$ is easy for some $a_0$.
  \item \emph{Differentiate in $a$} (\cref{prop:diffint}). The family is
        chosen so that $F'$ is easier than $F$: either $F'$ can be
        computed directly, or $F'$ is expressed through $F$ itself.
  \item \emph{Go back from $F'$ to $F$}: integrate $F'$ from $a_0$, or
        solve the differential equation that links $F'$ to $F$, with the
        known value $F(a_0)$ as initial condition. Then evaluate at
        $a_1$.
\end{enumerate}
All the inventiveness is in steps 1 and 2: where to put the parameter.
The examples show the two ways step 3 can go.

\section{When the derivative is easy}
\label{sec:feyn-easy}

\begin{example}{A logarithm in the denominator}{lntwo}
Compute $I=\int_0^1\frac{x-1}{\ln x}\,\mathrm dx$.

The $\ln x$ in the denominator is the obstacle. Differentiating $x^a$
with respect to $a$ produces a factor $\ln x$, so put the parameter in
the exponent:
\[
  F(a)=\int_0^1\frac{x^a-1}{\ln x}\,\mathrm dx ,\qquad a\in J=(-\tfrac12,2).
\]
Then $I=F(1)$, and $F(0)=0$ because the integrand is $0$ at $a=0$.
The derivative of the integrand is
$\partial_a\bigl((x^a-1)/\ln x\bigr)=x^a$: the obstacle cancels.
Hypotheses of \cref{prop:diffint}, on $S=(0,1)$: for $a\in J$,
$0<x^a\le x^{-1/2}$, which is integrable on $(0,1)$, so (3) holds with
$g(x)=x^{-1/2}$; and $(x^a-1)/\ln x=\int_0^ax^b\,\mathrm db$ is bounded
in absolute value by $\abs a\,x^{-1/2}$, so (1) holds. Hence
\[
  F'(a)=\int_0^1x^a\,\mathrm dx=\frac1{a+1},
  \qquad
  F(a)=F(0)+\int_0^a\frac{\mathrm db}{b+1}=\ln(1+a),
  \qquad
  I=F(1)=\ln2 .
\]
\end{example}

\begin{example}{The moments of a Gaussian}{gaussmoments}
Here the known member is a whole family. For $a>0$, the density of
$\Normal(0,1/(2a))$ integrates to $1$, that is
\[
  G(a)=\int_\R e^{-ax^2}\,\mathrm dx=\sqrt{\pi}\,a^{-1/2}.
\]
On an interval $(a_0,\infty)$ with $a_0>0$, every derivative
$\partial_a^k e^{-ax^2}=(-x^2)^ke^{-ax^2}$ is bounded by the integrable
$x^{2k}e^{-a_0x^2}$, so \cref{prop:diffint}, applied $k$ times, gives
\[
  \int_\R x^{2k}e^{-ax^2}\,\mathrm dx=(-1)^kG^{(k)}(a)
  =\sqrt\pi\;\frac{1\cdot3\cdots(2k-1)}{2^k}\;a^{-k-1/2}.
\]
Dividing by $G(a)$ and setting $a=1/(2\sigma^2)$: for
$\theta\sim\Normal(0,\sigma^2)$,
$\E[\theta^{2k}]=1\cdot3\cdots(2k-1)\,\sigma^{2k}$; in particular
$\E[\theta^2]=\sigma^2$ and $\E[\theta^4]=3\sigma^4$.
\end{example}

\section{When the derivative gives back the integral}
\label{sec:feyn-ode}

This is the case of Exercise~\ref{ex:ecos}. Let
$\theta\sim\Normal(0,\sigma^2)$, with density
$\phi(x)=e^{-x^2/(2\sigma^2)}/(\sigma\sqrt{2\pi})$, and let us look for
$\E[\cos\theta]$.

\emph{Why integration by parts alone does not close.} The density
satisfies $\phi'(x)=-(x/\sigma^2)\,\phi(x)$, so an integration by parts
turns $\E[\theta\sin\theta]$ into $\sigma^2\E[\cos\theta]$. That is one
equation between two unknowns. Integrating by parts again only brings in
$\E[\theta^2\cos\theta]$, then $\E[\theta^3\sin\theta]$, and so on.

\emph{Where to put the parameter.} The second unknown has the extra
factor $\theta$ in front of $\sin\theta$, and that factor is what
differentiating $\cos(a\theta)$ with respect to $a$ produces:
$\partial_a\cos(a\theta)=-\theta\sin(a\theta)$. So the parameter goes into
the frequency:
\[
  m(a)=\E[\cos(a\theta)]=\int_\R\cos(ax)\,\phi(x)\,\mathrm dx ,
  \qquad m(0)=1 .
\]
The second unknown becomes the derivative of the first, and the one
equation between two unknowns becomes a differential equation for one
unknown function.

\emph{Differentiate.} $\abs{\partial_a(\cos(ax)\phi(x))}
=\abs{x\sin(ax)}\,\phi(x)\le\abs x\,\phi(x)$, which is integrable: since
$x\,\phi(x)$ has the antiderivative $-\sigma^2\phi(x)$,
$\int_\R\abs x\,\phi(x)\,\mathrm dx=2\sigma^2\phi(0)=\sigma\sqrt{2/\pi}$. So
\cref{prop:diffint} applies on $J=\R$:
$m'(a)=-\int_\R x\sin(ax)\,\phi(x)\,\mathrm dx$.

\emph{Integrate by parts.} With $x\,\phi(x)=-\sigma^2\phi'(x)$, on an
interval $[-R_1,R_2]$,
\[
  \int_{-R_1}^{R_2}x\sin(ax)\,\phi(x)\,\mathrm dx
  =-\sigma^2\bigl[\sin(ax)\,\phi(x)\bigr]_{-R_1}^{R_2}
   +\sigma^2a\int_{-R_1}^{R_2}\cos(ax)\,\phi(x)\,\mathrm dx .
\]
As $R_1,R_2\to\infty$ the bracket tends to $0$, and both integrals tend
to the integrals over $\R$ by dominated convergence (bounds
$\abs x\,\phi(x)$ and $\phi(x)$). So
\[
  m'(a)=-\sigma^2a\,m(a).
\]

\emph{Solve.} $m(a)\,e^{\sigma^2a^2/2}$ has derivative
$\bigl(m'(a)+\sigma^2a\,m(a)\bigr)e^{\sigma^2a^2/2}=0$, so it is constant,
equal to $m(0)=1$:
\[
  m(a)=e^{-\sigma^2a^2/2},\qquad \E[\cos\theta]=m(1)=e^{-\sigma^2/2}.
\]
The trick gave more than was asked: the whole function $m$, not only
$m(1)$. In probability, $m$ is the real part of the
\emph{characteristic function} $a\mapsto\E[e^{\ii a\theta}]$ of
$\theta$, where $\ii$ is the imaginary unit ($\ii^2=-1$). The next
section computes the characteristic function directly.

\section{The route through complex numbers}
\label{sec:feyn-complex}

Euler's formula $e^{\ii y}=\cos y+\ii\sin y$ packs both expectations into
one:
\[
  \E[e^{\ii a\theta}]=\E[\cos(a\theta)]+\ii\,\E[\sin(a\theta)],
\]
where the integral of a complex-valued function is, by definition, the
integral of its real part plus $\ii$ times the integral of its imaginary
part. The exponential of a complex number $w$ satisfies
$e^{w_1+w_2}=e^{w_1}e^{w_2}$ and $\abs{e^w}=e^{\operatorname{Re}w}$.

\emph{Complete the square.} In the exponent of
$e^{\ii ax}\phi(x)$, with $c=a\sigma^2$,
\[
  \ii ax-\frac{x^2}{2\sigma^2}
  =-\frac{(x-\ii c)^2}{2\sigma^2}-\frac{a^2\sigma^2}2 ,
\]
as expanding $(x-\ii c)^2=x^2-2\ii cx-c^2$ shows. Hence
\[
  \E[e^{\ii a\theta}]
  =e^{-a^2\sigma^2/2}\;\frac{1}{\sigma\sqrt{2\pi}}
   \int_\R e^{-(x-\ii c)^2/(2\sigma^2)}\,\mathrm dx .
\]
If the integral were the one with $c=0$, it would be $\sigma\sqrt{2\pi}$
and the answer would be $e^{-a^2\sigma^2/2}$. The Gaussian has been
shifted by the imaginary amount $\ii c$; the lemma says that the shift
changes nothing.

\begin{lemma}{An imaginary shift of the Gaussian}{imagshift}
For every real $c$,
$\displaystyle H(c)=\int_\R e^{-(x-\ii c)^2/(2\sigma^2)}\,\mathrm dx
=\sigma\sqrt{2\pi}$.
\end{lemma}

\begin{proof}
The same trick, with $c$ as the parameter. Write
$k(c,x)=e^{-(x-\ii c)^2/(2\sigma^2)}$. Then
$\abs{k(c,x)}=e^{(c^2-x^2)/(2\sigma^2)}$, and
\[
  \partial_ck=\frac{\ii(x-\ii c)}{\sigma^2}\,k
  =-\ii\,\partial_xk .
\]
For $\abs c<C$, $\abs{\partial_ck}\le\frac{\abs x+C}{\sigma^2}\,
e^{(C^2-x^2)/(2\sigma^2)}$, which is integrable, so by
\cref{prop:diffint} (applied to real and imaginary parts),
\[
  H'(c)=-\ii\int_\R\partial_xk(c,x)\,\mathrm dx
  =-\ii\lim_{R_1,R_2\to\infty}\bigl[k(c,x)\bigr]_{x=-R_1}^{x=R_2}=0,
\]
where the integral over $\R$ is the limit of the integrals over
$[-R_1,R_2]$ by dominated convergence (same bound), and $k(c,\pm R)\to0$
because $\abs{k(c,x)}\to0$ as $\abs x\to\infty$. So $H$ is constant, and
$H(c)=H(0)=\sigma\sqrt{2\pi}$, the normalisation of the Gaussian density.
\end{proof}

Hence $\E[e^{\ii a\theta}]=e^{-a^2\sigma^2/2}$. Its real part is
$\E[\cos(a\theta)]=e^{-a^2\sigma^2/2}$, as in \cref{sec:feyn-ode}, and its
imaginary part is $\E[\sin(a\theta)]=0$, which recovers the symmetry
argument for the sine.

The complex route is shorter once the lemma is known, but the lemma is
the same trick in another form: the parameter is now the shift $c$, and
the derivative is a total derivative in $x$, whose integral is $0$. In
complex analysis the lemma is a case of Cauchy's theorem: the line of
integration can be moved from the real axis to the parallel line through
$\ii c$ without changing the integral. The proof above needs none of that
theory.

\section{Choosing the parameter}
\label{sec:feyn-choose}

There is no recipe, but the examples share a pattern.
\begin{itemize}[nosep]
  \item \emph{A known end.} The family must contain an easy member:
        $a=0$ in \cref{ex:lntwo} and in \cref{sec:feyn-ode}, $c=0$ in
        \cref{lem:imagshift}, all $a$ in \cref{ex:gaussmoments}.
  \item \emph{A derivative that removes the obstacle.} Look at what
        blocks the computation, and put $a$ where $\partial_a$ produces
        or cancels it: an exponent $x^a$ to cancel a $\ln x$, a
        frequency $\cos(ax)$ to produce a factor $x$, a scale
        $e^{-ax^2}$ to produce $x^2$. A damping factor $e^{-ax}$ is
        another common choice, to make an integral converge.
  \item \emph{A domination.} Check hypothesis (3) of
        \cref{prop:diffint} on an interval of parameters that contains
        both the known end and the target value.
  \item \emph{A way back.} Either integrate $F'$, or recognise a
        differential equation in $F$ and solve it from the known end.
\end{itemize}
